\chapter{IMPLEMENTAÇÃO E AMBIENTE DE DESENVOLVIMENTO}

A implementação do compilador MiniPar 2026.2 foi conduzida priorizando modularidade, desempenho e clareza didática, combinando a flexibilidade do ecossistema Python com o poder computacional de baixo nível do C++20.

\section{Tecnologias Utilizadas e Ambiente}

A Tabela \ref{tab:tecnologias} sumariza o conjunto de ferramentas, linguagens e frameworks empregados no projeto.

\begin{table}[htbp]
\centering
\caption{Tecnologias Empregadas na Construção do Compilador}
\label{tab:tecnologias}
\small
\begin{tabular}{p{3.2cm}p{3cm}p{7.5cm}}
\toprule
Tecnologia & Versão Utilizada & Finalidade no Projeto \\
\midrule
Python & 3.10 / 3.14 & Implementação do front-end, backend com otimizador, CLI, menu e servidor Web. \\
C++ & Padrão ISO C++20 & Linguagem-alvo transpilada com suporte nativo a threads e sockets. \\
GCC / MinGW & g++ 13+ & Compilação do código C++ gerado com flag de otimização \texttt{-O3}. \\
Flask & 3.0.0 & Servidor HTTP e API REST das duas interfaces Web. \\
D3.js & v7 (local) & Renderização gráfica vetorial interativa da Árvore Sintática (AST). \\
CodeMirror & v5 (local) & Editor de código no navegador com realce de sintaxe. \\
Pytest & 8.4.2 & Execução da suíte automatizada de testes de integração (117 casos). \\
Git / GitHub & 2.40+ & Controle de versão distribuído e colaboração da equipe. \\
\bottomrule
\end{tabular}
\end{table}

\subsection{Comentários Explicativos no Código-Fonte}
Conforme exigido na especificação, todo o código-fonte do compilador foi amplamente documentado. Em arquivos como \texttt{src/lexer.py}, \texttt{src/parser.py}, \texttt{src/semantic.py} e \texttt{src/optimizer.py}, cada método possui \textit{docstrings} formais e anotações explicativas em português, detalhando as regras gramaticais da BNF atendidas, os critérios de inferência de tipos e as otimizações aplicadas.

\section{Passo a Passo da Execução dos Programas de Teste}

O fluxo de compilação e execução segue uma esteira intuitiva, reproduzida tanto no terminal quanto na interface Web didática:

\begin{enumerate}
    \item Seleção do Código-Fonte: o usuário seleciona um dos programas de teste disponíveis ou digita o código no editor.
    \item Análise por Fases Didáticas: ao acionar o botão de análise, o backend executa o pipeline em memória, exibindo as tabelas de \textit{Tokens}, a \textit{Árvore Sintática AST} em D3.js, o diagnóstico \textit{Semântico} e a \textit{Tabela de Símbolos}.
    \item Inspeção de TAC e C++: o usuário pode alternar entre as abas para inspecionar o código intermediário linear e o código-fonte C++20 gerado.
    \item Compilação Nativa e Execução: o usuário insere os valores de entrada no console e aciona \textit{Executar}. O compilador invoca o \texttt{g++ -O3}, gera o binário executável e exibe a saída.
\end{enumerate}

\begin{figure}[htbp]
    \centering
    \framebox[\textwidth][c]{%
        \parbox{0.9\textwidth}{%
            \centering\vspace{1.5cm}
            \textcolor{red}{<Inserir prints do passo a passo das telas da execução dos testes na interface Web Didática>}\\
            \vspace{0.2cm}
            \textit{\small Nota: Capturar telas ilustrando: (a) Seleção do teste; (b) Inspeção da árvore AST; (c) Painel de execução com os resultados no console.}
            \vspace{1.5cm}
        }
    }
    \caption{Passo a Passo da Execução dos Programas de Teste na Interface Web Didática}
    \label{fig:passo_a_passo_web}
\end{figure}
